\chapter{Musculoskeletal Modeling Conventions}
\label{ch:msk-conventions}

\epigraph{%
  In biomechanics, the first source of error is not the mathematics
  but the coordinate system.
}{--- ISB Standardization Committee}

\section{Anatomical Reference Frames}

Biomechanical analysis requires a standardized language for describing body position and
motion. The International Society of Biomechanics (ISB) has published recommendations
for joint coordinate systems \cite{Wu2002, Wu2005} that serve as the international
standard.

\begin{center}
\begin{tikzpicture}[scale=1.0, >=Stealth]
  % Global frame
  \draw[thick, ->, red] (0, 0, 0) -- (2, 0, 0) node[right] {$\hat{\bm{x}}_G$ (anterior)};
  \draw[thick, ->, green] (0, 0, 0) -- (0, 2, 0) node[above] {$\hat{\bm{y}}_G$ (superior)};
  \draw[thick, ->, blue] (0, 0, 0) -- (0, 0, 2) node[below right] {$\hat{\bm{z}}_G$ (lateral)};

  \node at (-0.3, -0.3, 0) {$\{G\}$};

  % Segment frame (offset)
  \draw[thick, ->, red!60] (3, 1, 0) -- (4.5, 1.2, 0) node[right, font=\small] {$\hat{\bm{x}}_S$};
  \draw[thick, ->, green!60] (3, 1, 0) -- (3, 2.5, 0) node[above, font=\small] {$\hat{\bm{y}}_S$};
  \draw[thick, ->, blue!60] (3, 1, 0) -- (3, 1, 1.5) node[right, font=\small] {$\hat{\bm{z}}_S$};

  \node at (2.7, 0.5, 0) {$\{S\}$};

  % Bony landmarks on segment
  \filldraw[fill=black] (3, 1, 0) circle (3pt) node[below, font=\tiny] {Origin};
  \filldraw[fill=gray] (3.5, 2.8, 0) circle (2pt);
  \filldraw[fill=gray] (3.8, 2.5, 0) circle (2pt);

  \node[font=\tiny, gray] at (4.2, 2.9) {Landmarks};
\end{tikzpicture}
\end{center}

\subsection{The Anatomical Position}

All references to body position begin from the \textbf{anatomical position}: the subject
stands upright, feet together, arms at sides, palms facing forward. This defines the
\emph{zero configuration} from which all joint angles are measured.

\begin{definition}{Global Reference Frame}{global-frame}
  The global (laboratory) reference frame $\{G\}$ is defined as:
  \begin{itemize}
    \item $\hat{\bm{x}}_G$: anterior (forward)
    \item $\hat{\bm{y}}_G$: superior (upward)
    \item $\hat{\bm{z}}_G$: lateral right
  \end{itemize}
  This is a right-handed coordinate system. Note that different communities use
  different conventions: biomechanics typically uses $Y$-up, while robotics often
  uses $Z$-up. \textbf{Always verify the convention before combining data from
  different sources.}
\end{definition}

\subsection{Segment Reference Frames}

Each body segment has a local reference frame attached to it, defined by bony landmarks
that can be palpated or identified from medical imaging. The ISB recommendations specify
these landmarks for each segment \cite{Wu2002, Wu2005}.

\begin{example}{The Thigh Reference Frame}{thigh-frame}
  The ISB-recommended thigh reference frame $\{T\}$ is defined by:
  \begin{itemize}
    \item \textbf{Origin}: midpoint between medial and lateral femoral epicondyles
    \item $\hat{\bm{y}}_T$: from origin to hip joint center (proximal direction)
    \item $\hat{\bm{z}}_T$: perpendicular to the plane formed by the two epicondyles
          and the hip center, pointing laterally
    \item $\hat{\bm{x}}_T$: cross product $\hat{\bm{y}}_T \times \hat{\bm{z}}_T$
          (anterior direction)
  \end{itemize}
  The hip joint center is estimated from the pelvis landmarks using regression
  equations (e.g., \cite{Harrington2007}).
\end{example}

\subsection{Connection to Screw Theory}

In the language of Volume~0, each segment frame $\{S_i\}$ is related to the global
frame $\{G\}$ by a homogeneous transformation matrix $T_{GS_i} \in \SE$:
\begin{equation}
  T_{GS_i} = \begin{bmatrix} R_{GS_i} & \bm{p}_{GS_i} \\ \bm{0}^T & 1 \end{bmatrix},
  \label{eq:ch2:homogeneous}
\end{equation}
where $R_{GS_i} \in \SO$ is the rotation matrix and $\bm{p}_{GS_i} \in \Reals^3$ is
the position of the segment origin in the global frame.

The \emph{relative configuration} of segment $j$ with respect to proximal segment $i$
is:
\begin{equation}
  T_{S_i S_j} = T_{GS_i}^{-1} T_{GS_j}.
  \label{eq:ch2:relative-config}
\end{equation}
Joint angles are extracted from this relative transformation.

\section{Joint Coordinate Systems}

The ISB recommends decomposing the relative rotation $R_{S_i S_j}$ into three
angles using a specific Euler angle sequence for each joint. The choice of
sequence is \emph{not arbitrary}---it is selected to align with the clinical
definitions of flexion-extension, abduction-adduction, and internal-external
rotation.

\subsection{The Grood--Suntay Convention for the Knee}

The knee joint coordinate system, originally proposed by Grood and Suntay (1983)
\cite{GroodSuntay1983}, decomposes knee motion into:

\begin{enumerate}
  \item \textbf{Flexion-Extension} ($\alpha$): rotation about the femoral medio-lateral
        axis (the ``fixed'' axis of the femur)
  \item \textbf{Abduction-Adduction} ($\beta$): rotation about the ``floating'' axis
        (perpendicular to both the femoral and tibial body-fixed axes)
  \item \textbf{Internal-External Rotation} ($\gamma$): rotation about the tibial
        longitudinal axis (the ``body-fixed'' axis of the tibia)
\end{enumerate}

This is a ZXY Euler angle sequence when expressed in the Grood--Suntay convention.
The key insight is that the floating axis is neither fixed in the femur nor in the
tibia---it is defined by the cross product of the other two axes.

\begin{lstlisting}[language=Python, caption={Computing knee joint angles using the Grood-Suntay convention.}]
import numpy as np
from typing import Tuple

def grood_suntay_knee_angles(
    R_femur: np.ndarray,
    R_tibia: np.ndarray
) -> Tuple[float, float, float]:
    """Compute knee joint angles using Grood-Suntay convention.

    Parameters
    ----------
    R_femur : np.ndarray, shape (3, 3)
        Rotation matrix of femur segment frame in global coordinates.
    R_tibia : np.ndarray, shape (3, 3)
        Rotation matrix of tibia segment frame in global coordinates.

    Returns
    -------
    Tuple[float, float, float]
        (flexion_deg, adduction_deg, internal_rotation_deg)

    References
    ----------
    Grood, E.S. and Suntay, W.J. (1983). A joint coordinate system
    for the clinical description of three-dimensional motions.
    J. Biomechanical Engineering, 105(2), 136-144.
    """
    # Femoral medio-lateral axis (e1) - flexion axis
    e1 = R_femur[:, 2]  # lateral axis of femur

    # Tibial longitudinal axis (e3) - rotation axis
    e3 = R_tibia[:, 1]  # proximal axis of tibia

    # Floating axis (e2) - perpendicular to both
    e2 = np.cross(e3, e1)
    e2 = e2 / np.linalg.norm(e2)

    # Flexion-extension angle
    # Angle between femur anterior axis and floating axis
    femur_anterior = R_femur[:, 0]
    flexion = np.arctan2(
        np.dot(e2, R_femur[:, 1]),
        np.dot(e2, femur_anterior)
    )

    # Abduction-adduction angle
    # Angle between e1 and e3 minus 90 degrees
    sin_abd = -np.dot(e1, e3)
    cos_abd = np.linalg.norm(np.cross(e1, e3))
    adduction = np.arcsin(np.clip(sin_abd, -1.0, 1.0))

    # Internal-external rotation
    # Angle between floating axis and tibia anterior axis
    tibia_anterior = R_tibia[:, 0]
    internal_rot = np.arctan2(
        np.dot(e2, R_tibia[:, 2]),
        np.dot(e2, tibia_anterior)
    )

    return (
        np.degrees(flexion),
        np.degrees(adduction),
        np.degrees(internal_rot)
    )
\end{lstlisting}

\subsection{Multi-Representation Joint Angles}

Following the multi-representation philosophy of this series, we express joint
rotations in all four standard representations:

\textbf{1. Euler Angles (ISB Convention):}
\begin{equation}
  R_{ij} = R_z(\alpha) \, R_x(\beta) \, R_y(\gamma),
  \label{eq:ch2:euler-isb}
\end{equation}
where $\alpha, \beta, \gamma$ are the clinical angles (flexion, adduction, rotation).

\textbf{2. Rotation Matrix (SO(3)):}
\begin{equation}
  R_{ij} = \exp([\bm{\omega}]_\times \theta) \in \SO,
  \label{eq:ch2:rotmat}
\end{equation}
which avoids gimbal lock but requires 9 parameters with 6 constraints.

\textbf{3. Unit Quaternion:}
\begin{equation}
  \bm{q} = (\cos\tfrac{\theta}{2},\; \bm{\omega}\sin\tfrac{\theta}{2}),
  \quad \norm{\bm{q}} = 1,
  \label{eq:ch2:quaternion}
\end{equation}
which is singularity-free and computationally efficient.

\textbf{4. Screw Axis (Primary):}
\begin{equation}
  T_{ij} = e^{[\screw]\theta} \in \SE,
  \quad \screw = (\bm{\omega}, \bm{v}) \in \Reals^6,
  \label{eq:ch2:screw}
\end{equation}
which captures the coupled rotation-translation of biological joints naturally.

\begin{comparison}{Why Screw Axis Theory Is Natural for Biological Joints}{screw-bio}
  Unlike engineering revolute joints (pure rotation about a fixed axis), biological
  joints exhibit \emph{coupled} rotation and translation. The knee slides anteriorly
  as it flexes; the glenohumeral joint translates inferiorly during abduction.
  Screw axis theory captures this coupling naturally: every rigid-body motion is a
  rotation about and translation along a screw axis. The instantaneous screw axis
  (ISA) of a biological joint moves through space as the joint moves, unlike the
  fixed axis of an ideal revolute.

  This is precisely why Park and Lynch's \emph{Modern Robotics} framework is
  preferred: it handles coupled rotation-translation without special cases.
\end{comparison}

\section{Body Segment Parameters}

\subsection{Regression Equations}

The most widely used body segment parameter (BSP) data comes from:
\begin{enumerate}
  \item \textbf{Dempster (1955)}: cadaver study of 8 elderly males. Still widely cited
        but limited demographic range.
  \item \textbf{Zatsiorsky \& Seluyanov (1983, 1990)}: gamma-ray scanning of 100 living
        young males. More representative but uses different segment definitions.
  \item \textbf{de Leva (1996)}: adjustment of Zatsiorsky data to standard segment
        definitions. The current gold standard for adult male and female BSP.
  \item \textbf{Pavol et al.\ (2002)}: elderly adult data.
  \item \textbf{Jensen (1986, 1989)}: pediatric data.
\end{enumerate}

\subsection{Computing the Mass Matrix}

Given segment inertial properties, the mass matrix $M(\bm{q}) \in \Reals^{n \times n}$
of the multi-body system is computed using the composite rigid-body algorithm
(Volume~0, Chapter~10). In the body-fixed frame of segment $i$, the spatial
inertia is:
\begin{equation}
  \mathcal{G}_i = \begin{bmatrix} I_i & m_i [\bm{c}_i]_\times^T \\
  m_i [\bm{c}_i]_\times & m_i \mathbb{I}_3 \end{bmatrix} \in \Reals^{6 \times 6},
  \label{eq:ch2:spatial-inertia}
\end{equation}
where $I_i$ is the $3 \times 3$ rotational inertia about the segment frame origin,
$m_i$ is the segment mass, $\bm{c}_i$ is the center of mass position in the segment
frame, and $[\cdot]_\times$ denotes the skew-symmetric matrix.

\section*{Exercises}

\begin{enumerate}
  \item \textbf{Coordinate System Verification.}
        Verify that the Grood--Suntay knee joint coordinate system produces zero
        angles when the femur and tibia reference frames are aligned. Test with
        known rotations.

  \item \textbf{Euler Angle Singularity.}
        For the ZXY Euler angle sequence used in the knee, identify the gimbal lock
        configuration. Is this configuration ever reached during normal gait?
        During deep squatting?

  \item \textbf{Multi-Representation Conversion.}
        Write Python functions to convert between all four representations
        (Euler ZXY, rotation matrix, quaternion, screw axis) for a knee joint angle
        of $60^\circ$ flexion, $5^\circ$ adduction, $10^\circ$ internal rotation.

  \item \textbf{BSP Sensitivity.}
        Using the de~Leva segment properties, compute the knee extension torque
        required to hold the lower leg horizontal (isometric). Repeat using the
        Dempster data. What is the percentage difference? Discuss the implications for
        inverse dynamics accuracy.

  \item \textbf{(Programming.)} Implement the spatial inertia matrix computation
        for a 3-segment leg model (thigh, shank, foot). Verify positive-definiteness.

  \item \textbf{(Programming.)} Using motion capture data (or simulated data),
        compute knee joint angles from segment orientations using the Grood--Suntay
        code provided. Plot flexion angle over one gait cycle.
\end{enumerate}
